\section{Related work}\label{sec:related}

\paragraph{Evaluation pitfalls in ML-based NIDS} Sommer and Paxson~\cite{sommer2010outside} argued early that anomaly detection in networks differs fundamentally from typical ML settings: errors are expensive, the base rate of attacks is low~\cite{axelsson2000base}, and there is a semantic gap between alerts and actionable incidents. Arp et al.~\cite{arp2022dos} systematised pitfalls such as sampling bias, data snooping and inappropriate baselines, and Apruzzese et al.~\cite{apruzzese2023sok} called for assessments that reflect operational constraints. Several works document quality problems in the benchmarks themselves: labelling and feature-extraction errors in CIC-IDS2017~\cite{engelen2021troubleshooting,lanvin2023errors}, and ``design smells'' such as poor diversity and ill-defined generalisation criteria across seven datasets~\cite{flood2024bad}. Standardised NetFlow feature sets were proposed to enable comparison across datasets~\cite{sarhan2021netflow,sarhan2022standard,sarhan2022evaluating}, and temporal re-analyses of those datasets have recently been released~\cite{luay2025temporal}. In malware detection, TESSERACT~\cite{pendlebury2019tesseract} showed how spatial and temporal bias inflates results, and Apruzzese et al.~\cite{apruzzese2024adversarial} examined how drift interacts with adversarial perturbation in NIDS.

\paragraph{Cross-dataset generalisation} Layeghy and Portmann~\cite{layeghy2023explainable} and Cantone et al.~\cite{cantone2024cross} report near-perfect in-dataset performance that drops to near chance across datasets. Xin and Xu~\cite{xin2025crossdataset} add Platt-scaled calibration and AUC-oriented training to a transformer and evaluate across NSL-KDD, UNSW-NB15 and CIC-IDS2017. Jin and Lee~\cite{jin2026benchmark} harmonise features across HIKARI-2021, CIC-IDS2017 and CICIoT2023 and show that a 5\% false-positive-rate threshold chosen in-domain does not hold after transfer, even when ranking metrics look good. These studies treat detection as binary and do not examine unseen attack families or selective or conformal alerting.

\paragraph{Unseen attacks and uncertainty} Open-set NIDS approaches reject unknown attacks with clustering~\cite{farrukh2023detecting} or prototype learning with extreme-value thresholds~\cite{yu2025industrial}. Talpini et al.~\cite{talpini2024enhancing} compare uncertainty-quantification methods, including Bayesian neural networks, for detecting held-out attack classes, and Pal et al.~\cite{pal2025calibrated} study calibration and uncertainty for IoT attack detection under within-dataset shift. General-purpose baselines for misclassification and out-of-distribution detection include maximum softmax probability~\cite{hendrycks2017baseline} and energy scores~\cite{liu2020energy}; Ovadia et al.~\cite{ovadia2019trust} showed that calibration degrades under dataset shift, with ensembles~\cite{lakshminarayanan2017simple} among the most robust options. Post-hoc calibration methods include Platt scaling~\cite{platt1999probabilistic}, isotonic regression~\cite{zadrozny2002transforming,niculescu2005predicting} and Bayesian binning~\cite{naeini2015obtaining}.

\paragraph{Conformal and selective prediction in security} Conformal prediction provides distribution-free guarantees under exchangeability~\cite{vovk2022algorithmic,lei2018distribution,angelopoulos2023conformal}, with extensions to covariate shift~\cite{tibshirani2019conformal}, online adaptation~\cite{gibbs2021adaptive}, non-exchangeable data~\cite{barber2023conformal}, risk control~\cite{angelopoulos2024conformal} and outlier testing with conformal p-values~\cite{bates2023testing}. Selective classification trades coverage for accuracy~\cite{geifman2017selective}. In security, Transcend~\cite{jordaney2017transcend,barbero2022transcending} and CADE~\cite{yang2021cade} use conformal evaluators or contrastive embeddings to detect drifting malware samples. For intrusion detection, Yumlembam et al.~\cite{yumlembam2024comprehensive} add a conformal rejection layer to an adversarially hardened botnet detector, and Li~\cite{li2026robust} restores conformal coverage under adversarial feature perturbation for LLM-based detectors. Neither considers natural cross-dataset shift, unseen families or an analyst budget.

\paragraph{Positioning} Table~\ref{tab:positioning} summarises the dimensions covered by the closest works. To our knowledge, this is the first study to evaluate cross-dataset shift, unseen attack families, calibration and conformal or selective alerting \emph{jointly} and under an explicit analyst budget, as called for by a recent review~\cite{lew2026deployment}.

\begin{table}[t]
\centering
\caption{Coverage of deployment concerns in closely related work (\checkmark: evaluated; $\circ$: partially). XD: cross-dataset shift; OS: unseen attack families; CAL: probability calibration; AB: alert budget or fixed false-positive-rate operation; CP/SEL: conformal or selective prediction.}\label{tab:positioning}
\small
\begin{adjustbox}{max width=\linewidth}
\begin{tabular}{lccccc}
\toprule
Work & XD & OS & CAL & AB & CP/SEL\\
\midrule
Talpini et al.~\cite{talpini2024enhancing} & & \checkmark & \checkmark & & \\
Pal et al.~\cite{pal2025calibrated} & $\circ$ & $\circ$ & \checkmark & & \\
Xin and Xu~\cite{xin2025crossdataset} & \checkmark & & \checkmark & & \\
Cantone et al.~\cite{cantone2024cross} & \checkmark & & & & \\
Yumlembam et al.~\cite{yumlembam2024comprehensive} & & & & & \checkmark\\
Li~\cite{li2026robust} & $\circ$ & & & & \checkmark\\
Jin and Lee~\cite{jin2026benchmark} & \checkmark & & \checkmark & \checkmark & \\
\textbf{This work} & \checkmark & \checkmark & \checkmark & \checkmark & \checkmark\\
\bottomrule
\end{tabular}
\end{adjustbox}
\end{table}
